%-------------------------
% LaTeX-тегі түйіндеме
% Авторы : Рахим Баймурзин үшін бейімделген
% Құрастыру: tectonic cv-kk.tex   (немесе)  xelatex cv-kk.tex
%            Unicode қолдайтын қозғалтқыш қажет (XeTeX/LuaTeX).
%------------------------

\documentclass[letterpaper,11pt]{article}

\usepackage[empty]{fullpage}
\usepackage{titlesec}
\usepackage[usenames,dvipsnames]{color}
\usepackage{enumitem}
\usepackage[hidelinks]{hyperref}
\usepackage{fancyhdr}
\usepackage{comment}
\usepackage{array}
\usepackage{iftex}

% Computer Modern Unicode — шрифт суреті сол күйде, бірақ кириллицасы бар.
\usepackage{fontspec}
\setmainfont{cmunrm.otf}[
  BoldFont          = cmunbx.otf,
  ItalicFont        = cmunti.otf,
  BoldItalicFont    = cmunbi.otf,
  SmallCapsFont     = cmunrm.otf,
  SmallCapsFeatures = {Letters=SmallCaps},
]

% Қазақ тілінің тасымал үлгілері жоқ, сондықтан орыс тілінің үлгілері қолданылады.
\usepackage{polyglossia}
\setmainlanguage{russian}
\setotherlanguage{english}

% Тегтелген PDF глиф картасы: тек pdfTeX үшін, XeTeX/LuaTeX-те өткізіліп жіберіледі.
\ifPDFTeX
  \input{glyphtounicode}
  \pdfgentounicode=1
\fi

\pagestyle{fancy}
\fancyhf{}
\fancyfoot{}
\renewcommand{\headrulewidth}{0pt}
\renewcommand{\footrulewidth}{0pt}

\addtolength{\oddsidemargin}{-0.5in}
\addtolength{\evensidemargin}{-0.5in}
\addtolength{\textwidth}{1in}
\addtolength{\topmargin}{-.6in}
\addtolength{\textheight}{1.3in}

\urlstyle{same}
\raggedbottom
\raggedright
\setlength{\tabcolsep}{0in}

\titleformat{\section}{
  \vspace{-4pt}\scshape\raggedright\large
}{}{0em}{}[\color{black}\titlerule \vspace{-5pt}]

%-------------------------
% Түйіндеме макростары
%-------------------------

\newcommand{\resumeItem}[1]{
  \item\small{
    {#1 \vspace{-2pt}}
  }
}

\newcommand{\resumeSubheading}[4]{
  \vspace{-2pt}\item
    \begin{tabular*}{0.97\textwidth}[t]{l@{\extracolsep{\fill}}r}
      \textbf{#1} & #2 \\
      \textit{\small#3} & \textit{\small #4} \\
    \end{tabular*}\vspace{-7pt}
}

\newcommand{\resumeProjectHeading}[2]{
    \item
    \begin{tabular*}{0.97\textwidth}{l@{\extracolsep{\fill}}r}
      \small#1 & #2 \\
    \end{tabular*}\vspace{-7pt}
}

\renewcommand\labelitemii{$\vcenter{\hbox{\tiny$\bullet$}}$}


\newcommand{\resumeSubHeadingListStart}{\begin{itemize}[leftmargin=0.15in, label={}]}
\newcommand{\resumeSubHeadingListEnd}{\end{itemize}}
\newcommand{\resumeItemListStart}{\begin{itemize}}
\newcommand{\resumeItemListEnd}{\end{itemize}\vspace{-5pt}}

%-------------------------------------------
%%%%%%  ТҮЙІНДЕМЕ ОСЫ ЖЕРДЕН БАСТАЛАДЫ  %%%%%%%%%%%%%%%%%%%%%%%%%%%%

\begin{document}

\begin{center}
    \textbf{\Huge \scshape Рахим Баймурзин} \\
    Қазақстан $|$
    \href{mailto:baimurzzin@gmail.com}{\underline{baimurzzin@gmail.com}} $|$
    \href{https://www.linkedin.com/in/baimurzzin/}{\underline{linkedin}} $|$
    tg: @baimurzzin
\end{center}

%-----------БІЛІМ-----------
\section{Білім}
\resumeSubHeadingListStart
  % \resumeSubheading
  %   {Назарбаев Зияткерлік мектебі (НЗМ)}{Қарағанды, Қазақстан}
  %   {Орта білім аттестаты, GPA 3.75, ҰБТ 138/140, мектеп тарихындағы алғашқы ISO медалисі}{2016 -- 2022}
  \resumeSubheading
    {Назарбаев Университеті}{Астана, Қазақстан}
    {Қолданбалы математика магистрі}{2026 -- 2028}
  \resumeSubheading
    {Назарбаев Университеті}{Астана, Қазақстан}
    {Математика бакалавры, GPA 3.21}{2022 -- 2026}
\resumeSubHeadingListEnd

%-----------МАРАПАТТАР-----------
\section{Марапаттар мен жетістіктер}
\resumeSubHeadingListStart
  \item{\begin{tabular}{@{}>{\raggedright\arraybackslash}p{0.45\textwidth}@{\hspace{0.05\textwidth}}>{\raggedright\arraybackslash}p{0.45\textwidth}@{}}
    \textbf{IMO 2021 -- қола}                       & \textbf{ICPC (NEF), аймақтық финалист}                \\
    IMC 2026 -- күміс                               & Республикалық студенттік МО 2026 (1/50)       \\
    MBZUAI K2 Think Хакатон, жеңімпаз (1/900)   & KAIST E5 Global Стартап Көрмесі (2/40)
  \end{tabular}}
\resumeSubHeadingListEnd

%-----------ЖҰМЫС ТӘЖІРИБЕСІ-----------
\section{Жұмыс тәжірибесі}
\resumeSubHeadingListStart

  \resumeSubheading
    {Desargues AI}{Астана, Қазақстан}
    {Құрылтайшы}{ақп 2026 -- қазір}
    \resumeItemListStart
      \resumeItem{Жасанды математикалық суперинтеллект жасап жатырмын.}
    \resumeItemListEnd

  \resumeSubheading
    {Higgsfield AI}{Алматы, Қазақстан}
    {ML инженері}{қыр 2025 -- ақп 2026}
    \resumeItemListStart
      \resumeItem{Датасеттер жинап, ЖИ жасаған және нақты кескіндерді ажырататын, recall көрсеткіші жоғары өнімдік классификаторлар жасадым.}
    \resumeItemListEnd

  \resumeSubheading
    {Mercor AI Lab}{Қашықтан}
    {IMO және математикалық зерттеулер сарапшысы / дерек авторы}{2025}
    \resumeItemListStart
      \resumeItem{FrontierMath рейтингіндегі ең мықты модельді құру үшін деректер жинап, олардың математикалық сапасын бағаладым.}
    \resumeItemListEnd

  \resumeSubheading
    {Математика олимпиадалары}{Қазақстан}
    {Математика жаттықтырушысы, жиын ұйымдастырушысы, қазы, есеп авторы}{2021 -- қазір}
    \resumeItemListStart
      \resumeItem{Ұлттық IMO құрамасы үшін бірнеше жаттығу жиынын өткіздім.}
      \resumeItem{4 оқушыны бастапқы деңгейден халықаралық медаль деңгейіне дейін дайындадым.}
      \resumeItem{20-дан астам аймақтық, республикалық және халықаралық олимпиадада қазы әрі ұйымдастырушы болдым.}
    \resumeItemListEnd

\resumeSubHeadingListEnd

%-----------АКАДЕМИЯЛЫҚ ТӘЖІРИБЕ-----------
\section{Академиялық тәжірибе}
\resumeSubHeadingListStart

  \resumeSubheading
    {Назарбаев Университеті}{Астана, Қазақстан}
    {Микроэкономика бойынша оқытушы көмекшісі}{күз 2024}
    \resumeItemListStart
      \resumeItem{Семинарлар жүргізіп, студенттердің емтиханға дайындалуына көмектестім.}
    \resumeItemListEnd

  \vspace{-18pt}

  \resumeSubheading
    {}{}
    {Математика бойынша ғылыми көмекші}{күз 2025}
    \resumeItemListStart
      \resumeItem{Рустем Таханов және Жеңісбек Асылбеков профессорлармен бірге үлестірімдердің келісім критерийлері бойынша жұмыс істедім. Критерийлердің қуатын өлшеу үшін таңдамалы үлестірімдерді құрадым.}
    \resumeItemListEnd

  \vspace{-18pt}

  \resumeSubheading
    {}{}
    {Бакалавриат зерттеушісі}{көктем 2026}
    \resumeItemListStart
      \resumeItem{Профессор Адильбек Каиржанның жетекшілігімен дисперсиялық дербес туындылы теңдеулер бойынша дипломдық жұмыс жаздым.}
    \resumeItemListEnd

  \vspace{-18pt}

  \resumeSubheading
    {}{}
    {15-ші ISAAC конгресінің еріктісі}{тамыз 2025}
    \resumeItemListStart
      \resumeItem{Шетелдік қонақтарды алып жүрдім, сұхбаттар жүргіздім, семинарлар ұйымдастырып, оларға қатыстым.}
    \resumeItemListEnd

  \resumeSubheading
    {Imperial College London}{Лондон, Ұлыбритания}
    {Математиканы формализациялау сарапшысы (LEAN)}{шілде 2026}
    \resumeItemListStart
      \resumeItem{Кевин Баззард өткізген Ферманың ұлы теоремасын формализациялау практикумына қатыстым. Репрезентациялар теориясын формализацияладым.}
    \resumeItemListEnd

\resumeSubHeadingListEnd

% %-----------ЕРІКТІЛІК ТӘЖІРИБЕСІ-----------
% \section{Еріктілік тәжірибесі}
% \resumeSubHeadingListStart
%
%   \resumeSubheading
%     {NU SIAM Student Chapter}{Астана, Қазақстан}
%     {Мүше және математика өкілі}{2022 -- 2026}
%     \resumeItemListStart
%       \resumeItem{Математикалық баяндамалар, жарыстар (мысалы, Integration Bee) ұйымдастырдым және факультет талқылауларында математика студенттерінің атынан сөйледім.}
%       \resumeItem{1000-нан астам жазылушысы бар академиялық Telegram арнасын құрып, жүргіздім.}
%     \resumeItemListEnd
%
% \resumeSubHeadingListEnd

%-----------КУРСТАР-----------
% \section{Тиісті курстар}
% \resumeSubHeadingListStart
%   \item{\begin{tabular}{@{}p{0.48\textwidth} p{0.48\textwidth}@{}}
%     MATH 465: дифференциалдық геометрия & MATH 481: дербес туындылы теңдеулер \\
%     MATH 510: өлшемдер теориясы          & MATH 602: аналитикалық сандар теориясы \\
%     MATH 682: функционалдық талдау       & CSCI 371: алгоритмдер II              \\
%     CSCI 694: терең оқыту                & ECON 498: жалпы тепе-теңдік           \\
%     CHEM 211: органикалық химия          & MINE 402: көмір өндіру
%   \end{tabular}}
% \resumeSubHeadingListEnd

\begin{comment}
%-----------ЖОБАЛАР-----------
\section{Жобалар}
\resumeSubHeadingListStart

  \resumeProjectHeading
    {\textbf{Тері обырын анықтау}}{\textit{күз 2025}}
    \resumeItemListStart
      \resumeItem{Медициналық кескіндер бойынша тері обырын анықтайтын CNN негізіндегі модель жасадым.}
      \resumeItem{Терең оқытудың медициналық диагностикадағы қолданысын зерттедім.}
    \resumeItemListEnd

\resumeSubHeadingListEnd
\end{comment}

%-----------ДАҒДЫЛАР МЕН ҚЫЗЫҒУШЫЛЫҚТАР-----------
\section{Басқа}
\resumeSubHeadingListStart
  \resumeItem{\textbf{Дағдылар:} ML, DS, GitHub, C++, Python, AWS, Docker, SQL, BigQuery, Grafana, LaTeX, Excel, Apify.}
  % \resumeItem{\textbf{Қызығушылықтар:} спорттық бағдарламалау (Codeforces 1600), шахмат (1700 блиц/рапид), фортепиано, үстел теннисі.}
\resumeSubHeadingListEnd

\end{document}
